\documentclass[../../main.tex]{subfiles}
\begin{document}

{\bf[解]}

\paragraph{(1)}

% TODO: upgrade
\begin{figure}[htb]
  \centering
  \begin{tikzpicture}[scale=1.2, >=stealth, every node/.style={font=\small}]

    % --- 1. サイコロの目（1〜6）の垂直軸 ---
    \draw[->, thick, gray!80!black] (0, 0.4) -- (0, 6.7) node[above] {サイコロの出目};
    \foreach \y in {1, 2, 3, 4, 5, 6} {
      \draw[gray!40] (-0.1, \y) -- (7.2, \y);
      \node[left=2pt, font=\scriptsize] at (0, \y) {$\y$};
    }

    % 出目 k の基準ライン (例として k = 4)
    \def\kVal{4}
    \draw[dashed, red!85!black, thick] (0, \kVal) -- (7.2, \kVal)
    node[right=2pt, font=\footnotesize, red!90!black] {出目 $k$};

    % --- 2. 許容範囲（<= k）のハイライト ---
    \fill[blue!8, opacity=0.8] (0.5, 0.5) rectangle (5.3, \kVal);

    % --- 3. 試行回数のステップ (1 回目 〜 n 回目) ---
    % x 座標: 1回目(1.0), 2回目(2.2), ... , n-1回目(4.8), n回目(6.3)
    \coordinate (T1) at (1.0, 0);
    \coordinate (T2) at (2.2, 0);
    \coordinate (Tm) at (3.5, 0); % 省略
    \coordinate (Tn1) at (4.8, 0);
    \coordinate (Tn) at (6.3, 0);

    % ステップの目盛りラベル
    \node[below=3pt, font=\footnotesize] at (T1) {$1$ 回目};
    \node[below=3pt, font=\footnotesize] at (T2) {$2$ 回目};
    \node[below=3pt, font=\normalsize] at (Tm) {$\dots$};
    \node[below=3pt, font=\footnotesize] at (Tn1) {$n-1$ 回目};
    \node[below=3pt, font=\footnotesize\bfseries, red!85!black] at (Tn) {$n$ 回目};

    % --- 4. 各回の出目の制約プロット ---
    % 1回目〜n-1回目: すべて k 以下
    \foreach \x/\val in {1.0/2, 2.2/4, 4.8/3} {
      \draw[line width=1.5pt, teal!85!black, {Bar[width=4pt]}-{Bar[width=4pt]}]
      (\x, 0.8) -- (\x, \kVal);
      \fill[teal!85!black] (\x, \val) circle (2.0pt);
    }
    \node[teal!90!black, font=\scriptsize, align=center] at (2.9, 2.2)
    {各回とも\\$1 \sim k$ の目が出現\\[2pt]\textbf{確率 $\frac{k}{6}$}};

    % n 回目: ちょうど目 k が出る
    \fill[red!85!black] (6.3, \kVal) circle (2.5pt);
    \draw[->, thick, red!85!black] (6.3, \kVal + 0.9) -- (6.3, \kVal + 0.15)
    node[midway, right=1pt, font=\scriptsize] {ちょうど目 $k$ が出る};
    \node[above=8pt, font=\footnotesize\bfseries, red!85!black] at (6.3, \kVal)
    {確率 $\frac{1}{6}$};

    % --- 5. 上部の確率積ブレース ---
    % 1回目〜n-1回目の (k/6)^{n-1}
    \draw[decorate, decoration={brace, amplitude=5pt}, thick, teal!85!black]
    (0.6, 6.2) -- (5.2, 6.2)
    node[midway, above=6pt, font=\footnotesize, teal!90!black] {
      $n-1$ 回連続で $k$ 以下 $\implies \left(\frac{k}{6}\right)^{n-1}$
    };

    % n回目の 1/6
    \draw[decorate, decoration={brace, amplitude=4pt}, thick, red!85!black]
    (5.8, 6.2) -- (6.8, 6.2)
    node[midway, above=6pt, font=\footnotesize, red!90!black] {
      $\frac{1}{6}$
    };

  \end{tikzpicture}
  \caption{事象 $a_n = k$ が起こる出目の推移}
  \label{fig:dice_stopping_time_model}
\end{figure}

$k$を$1$から$6$までの整数とする．$a_n=k$となる確率は$n=1$のとき$\dfrac{1}{6}$
である．$n\le 2$のときは$a_1$から$a_{n-1}$がすべて$k$以下となる時で
\begin{align}
  P(a_n=k)=\left(\dfrac{k}{6}\right)^{n-1}\dfrac{1}{6} \quad (n \ge 2)\label{eq:1}
\end{align}
である．これは$n=1$でも成立する．よって求める期待値の極限は
\begin{align}
  E(n)
  &=\sum_{k=1}^6kP(a_n=k) \\
  &=\sum_{k=1}^6\left(\dfrac{k}{6}\right)^n \\
  &\xrightarrow{n\to\infty} 1 \cdots\text{(答)}
\end{align}
となる．

\paragraph{(2)}

期待値の加法定理および\cref{eq:1}から
\begin{align}
  N(n)&=\sum_{l=1}^nP(a_l=2)　\\
  &=\sum_{l=1}^n\left(\dfrac{2}{6}\right)^{n-1}\dfrac{1}{6} \\
  &=\dfrac{1}{6}\dfrac{1-(1/3)^n}{1-1/3} \\
  &\xrightarrow{n\to\infty}\dfrac{1}{4}\cdots\text{(答)}
\end{align}
となる．

\newpage
\end{document}
